\textbf{Proof.}
We give explicit submodels.  Put \(f_S(x)=(a+1)x^a\) and \(f_T(x)=1\), so \(g_S=a+1\) and \(g_T=1\).  The inequalities imposed on the density constants make these factors admissible, and \(M>a+1\) puts both constant factors in their stated Hölder balls.  Retain the fixed randomized propensity \(e\).  Conditional on \(X=x\), let both source potential outcomes be Bernoulli, independent of \(A\), with
\[
 \mu_{0,\pm}(x)=\frac12,
 \qquad
 \mu_{1,\pm}(x)=\frac12\pm\delta b_h(x).
\]
Give the target potential outcomes the same conditional means.  Here \(\delta>0\) is fixed sufficiently small.  Scaling of a smooth bump gives \(\|b_h\|_{\mathcal H^s}\le C_K\), uniformly in \(h\).  Hence \(\delta\) can be chosen so that both armwise regressions lie in \(\mathcal H^s(M)\) and in \([0,1]\).  It can also be chosen to satisfy
\[
 \operatorname{Var}(Y\mid A=d,X=x)\ge
 \frac14-\delta^2\ge M^{-1};
\]
this is possible because \(M>4\).  Consistency, conditional randomization, mean transportability, and treatment positivity then hold by construction.  Thus both laws belong to the model class.

For \(|u|\le\delta<1/2\), direct differentiation of the Bernoulli log likelihood, or the inequality \(\operatorname{kl}(p,q)\le(p-q)^2/\{q(1-q)\}\), gives
\[
 \operatorname{kl}\!\left(\operatorname{Bern}\left(\frac12+u\right),
 \operatorname{Bern}\left(\frac12-u\right)\right)\le C_0u^2.
\]
The two source laws differ only for treated outcomes.  Tensorization and \(e(x)\le1\) therefore yield
\[
 \begin{aligned}
 \operatorname{KL}(P_{+,h}^{S,n},P_{-,h}^{S,n})
 &\le C_0n\delta^2\int_0^h b_h(x)^2(a+1)x^a\,dx\\
 &=C_0(a+1)n\delta^2h^{2s+a+1}
   \int_0^1K(u)^2u^a\,du
 \le Cnh^{2s+a+1}.
 \end{aligned}
\]
Since \(K\ge0\) is nonzero, \(k_1=\int_0^1K(u)\,du>0\), and consequently
\[
 |\tau_+-\tau_-|
 =2\delta\int_0^h b_h(x)\,dx
 =2\delta k_1h^{s+1}\ge ch^{s+1}.
\]

For the independent target-sample construction, let \(\ell(x)=x-1/2\), keep a fixed smooth nonconstant contrast \(\mu(x)=\delta_0\ell(x)\), and leave the source law unchanged.  Define
\[
 f_{T,\pm}(x)=1\pm u_m\ell(x),\qquad u_m=\delta_1m^{-1/2}.
\]
The densities integrate to one.  Because \(c_{T,-}<1<c_{T,+}\), a fixed sufficiently small \(\delta_1>0\) makes both densities positive, inside the prescribed pointwise bounds, and inside \(\mathcal H^t(M)\), for every \(m\).  The elementary bound for two densities bounded away from zero gives
\[
 \operatorname{KL}(P_{T,+}^{\otimes m},P_{T,-}^{\otimes m})
 \le C_1m u_m^2\int_0^1\ell(x)^2\,dx\le C_2,
\]
whereas
\[
 |\tau_+-\tau_-|
 =2\delta_0u_m\int_0^1\ell(x)^2\,dx
 \ge c m^{-1/2}.
\]
Assigning target potential outcomes with the fixed conditional means supplies the required full-data extensions. \(\square\)